\documentclass[11pt,a4paper]{article}

\usepackage[T1]{fontenc}
\usepackage[utf8]{inputenc}
\usepackage[margin=25mm]{geometry}
\usepackage{amsmath,amssymb}
\usepackage{booktabs}
\usepackage{siunitx}
\usepackage[expansion=false]{microtype}
\usepackage[hidelinks]{hyperref}

\newcommand{\thetae}{\theta_{e}}
\newcommand{\thetam}{\theta_{m}}

\title{Torque Ripple and Cogging Torque in Permanent-Magnet Synchronous Machines}
\author{}
\date{}

\begin{document}
\maketitle

\begin{abstract}
This note describes the physical origin of torque ripple in permanent-magnet
synchronous machines (PMSM), separates the current-independent contribution
(cogging torque) from the current-dependent contributions (back-EMF harmonics,
reluctance modulation, inverter and sensor imperfections), and derives the
harmonic orders associated with each mechanism. A dedicated section analyses how
the uncompensated ripple propagates through a back-to-back AFE\,+\,inverter
structure with a reduced DC-link capacitance, and quantifies the resulting
interharmonic emission on a \SI{50}{\hertz} grid for a machine running at a
fundamental frequency of \SI{15.6}{\hertz}. Design-side, control-side and
grid-side mitigation methods are summarised at the end.
\end{abstract}

\section{Torque production in the ideal machine}

In the ideal PMSM the phase flux linkages produced by the magnets are purely
sinusoidal functions of rotor position, the airgap permeance is constant, and
the machine is fed with sinusoidal balanced currents. Under these assumptions
the torque in the rotor reference frame is
\begin{equation}
T_{e} = \frac{3}{2}\,p\,\Big[\lambda_{pm}\,i_{q} + (L_{d}-L_{q})\,i_{d}i_{q}\Big],
\label{eq:ideal}
\end{equation}
where $p$ is the number of pole pairs, $\lambda_{pm}$ the magnet flux linkage
(peak, per phase), and $L_d$, $L_q$ the synchronous inductances. Both terms are
constant for constant $i_d$, $i_q$: the ideal machine produces no ripple.

Every ripple mechanism discussed below is therefore a violation of one of the
three assumptions: sinusoidal flux-linkage distribution, position-independent
permeance, or ideal current control.

\section{Energy-based formulation}

The complete instantaneous torque follows from the derivative of the magnetic
co-energy $W'(\mathbf{i},\thetam)$ at constant current:
\begin{equation}
T_{e}(\mathbf{i},\thetam) = \left.\frac{\partial W'(\mathbf{i},\thetam)}{\partial \thetam}\right|_{\mathbf{i}=\text{const}} .
\end{equation}
Writing $W'$ as the sum of a winding term, a mutual (magnet--winding) term and
a magnet-only term gives the standard decomposition
\begin{equation}
T_{e} = \underbrace{\mathbf{i}^{\mathsf{T}}\frac{\partial \boldsymbol{\lambda}_{pm}(\thetam)}{\partial \thetam}}_{\text{alignment (mutual) torque}}
\; + \;
\underbrace{\frac{1}{2}\,\mathbf{i}^{\mathsf{T}}\frac{\partial \mathbf{L}(\thetam,\mathbf{i})}{\partial \thetam}\,\mathbf{i}}_{\text{reluctance torque}}
\; + \;
\underbrace{\frac{\partial W'_{pm}(\thetam)}{\partial \thetam}}_{\text{cogging torque}} .
\label{eq:decomp}
\end{equation}

The three terms have fundamentally different behaviour:

\begin{itemize}
\item the first is proportional to current and reproduces the harmonic content
      of $\partial\boldsymbol{\lambda}_{pm}/\partial\thetam$, i.e.\ of the back-EMF;
\item the second is proportional to the square of the current and depends on
      the position (and current) dependence of the inductance matrix, i.e.\ on
      saliency, slotting and saturation;
\item the third is independent of the current and exists even with the machine
      unexcited. This is the cogging torque.
\end{itemize}

\section{Cogging torque}
\label{sec:cogging}

\subsection{Physical origin}

Cogging torque is the tendency of the rotor to align itself in positions of
minimum reluctance for the magnet flux. It exists because the stator is
slotted: as the rotor turns, the magnets see an airgap permeance that is
modulated by the alternation of teeth and slot openings, so the stored magnetic
energy of the magnet field varies with position. Any variation of energy with
position is a torque.

A first-order description uses the product of the magnet magnetomotive force
and the slotted-airgap permeance function referred to the stator:
\begin{equation}
F(\theta,\thetam)=\sum_{m}F_{m}\cos\!\big[mp(\theta-\thetam)\big],
\qquad
\Lambda(\theta)=\Lambda_{0}+\sum_{n\geq 1}\Lambda_{n}\cos\!\big(nQ\,\theta\big),
\end{equation}
with $Q$ the number of stator slots. The airgap energy contains products
$F_{m}\Lambda_{n}$; after integration over the airgap circumference, only the
terms satisfying
\begin{equation}
nQ = m\,p
\end{equation}
survive and depend on $\thetam$. The resulting torque is periodic with
\begin{equation}
N_{c} = \operatorname{lcm}(Q,\,2p)
\end{equation}
periods per mechanical revolution, and can be written as a Fourier series
\begin{equation}
T_{cog}(\thetam) = \sum_{k\geq 1} T_{cog,k}\,\sin\!\big(k N_{c}\thetam + \varphi_{k}\big).
\label{eq:cogging}
\end{equation}
The series has zero mean over one revolution: cogging torque does not
contribute average torque, only ripple.

Referred to the electrical period, the fundamental cogging order is
$N_{c}/p$. For $Q=12$, $2p=10$ ($p=5$): $N_{c}=60$ mechanical periods per
revolution, i.e.\ a $12^{\text{th}}$ electrical harmonic.

\subsection{Design factors}

The amplitude $T_{cog,1}$ decreases as $N_{c}$ increases, and $N_{c}$ is
maximised when $\gcd(Q,2p)$ is small; $\gcd(Q,2p)$ is therefore commonly used
as a cogging ``goodness'' indicator (lower is better). Other dominant factors
are the slot-opening-to-airgap ratio, the pole-arc to pole-pitch ratio, the
magnetisation profile, and skewing.

Continuous skew over an angle $\theta_{sk}$ attenuates the $k$-th cogging
harmonic by
\begin{equation}
k_{sk,k}=\frac{\sin\!\big(k N_{c}\theta_{sk}/2\big)}{k N_{c}\theta_{sk}/2},
\end{equation}
which vanishes for $\theta_{sk}=2\pi/N_{c}$, i.e.\ one full cogging period. Step
skew of the rotor magnets in $n_{s}$ axial blocks gives a comparable
attenuation with simpler manufacturing.

\section{Current-dependent torque ripple}
\label{sec:current}

\subsection{Back-EMF harmonics}

Real magnet flux-linkage distributions are not sinusoidal. Writing the phase
flux linkage as
\begin{equation}
\lambda_{a}(\thetae)=\sum_{h\ \text{odd}}\lambda_{h}\cos(h\thetae),
\end{equation}
the corresponding back-EMF is
$e_{a}=\omega_{e}\sum_{h}h\lambda_{h}\sin(h\thetae)$, i.e.\ the harmonic content
is amplified by the harmonic order $h$.

Feeding the machine with a purely sinusoidal current $i_{a}=I\sin(\thetae)$ (and
the corresponding $\pm 120^{\circ}$ phase shifts), the instantaneous torque is
obtained from the power balance $T_{e}=\big(e_{a}i_{a}+e_{b}i_{b}+e_{c}i_{c}\big)/\omega_{m}$.
Each product $\sin(h\thetae)\sin(\thetae)$ generates components of order $h-1$
and $h+1$; summing the three phases cancels every order that is not a multiple
of six, and the triplen flux-linkage harmonics produce only zero-sequence EMF,
which develops no torque in a star-connected machine without neutral. The result is
\begin{equation}
T_{e}(\thetae)=\frac{3}{2}\,p\,I\left[\lambda_{1}
+\sum_{n\geq 1}\Big((6n{+}1)\lambda_{6n+1}-(6n{-}1)\lambda_{6n-1}\Big)\cos(6n\thetae)\right].
\label{eq:6k}
\end{equation}

Equation \eqref{eq:6k} is the central result for this ripple family: the
$5^{\text{th}}$ and $7^{\text{th}}$ back-EMF harmonics produce a
$6^{\text{th}}$-order torque ripple, the $11^{\text{th}}$ and $13^{\text{th}}$
produce a $12^{\text{th}}$-order ripple, and so on. Note the sign: the two
contributions may partially cancel, so a machine with a distorted back-EMF is
not necessarily a machine with large ripple.

The dual statement also holds. If the back-EMF is sinusoidal but the current
contains $5^{\text{th}}$ and $7^{\text{th}}$ harmonics, the same
$6^{\text{th}}$-order ripple appears. In the rotor frame both cases appear as a
$6k$ ripple superimposed on $i_{d}$, $i_{q}$ and on the torque.

\subsection{Reluctance and saturation effects}

In interior-magnet machines the inductances are position- and
current-dependent because of slotting, rotor geometry and local saturation of
the bridges and tooth tips. Expanding $L_{d}(i_d,i_q,\thetae)$ and
$L_{q}(i_d,i_q,\thetae)$ in \eqref{eq:decomp} yields additional $6k$ terms whose
amplitude grows with the square of the current, so this contribution dominates
at high load while cogging dominates at light load. Cross-saturation
($\partial\lambda_{d}/\partial i_{q}\neq 0$) adds a further load-dependent
component and detunes the MTPA angle.

\subsection{Inverter and control-chain contributions}

These sources are independent of the machine design and are frequently the
dominant ones in a well-designed motor:

\begin{itemize}
\item \textbf{Dead time and device voltage drop.} The distortion of the applied
      voltage is a square-wave function of the current sign, generating
      $6k\pm1$ current harmonics and therefore $6k$ torque ripple. The effect
      is largest at low speed and low current, where the distortion voltage is
      comparable with the fundamental.
\item \textbf{Current-sensor DC offset.} An offset on one phase appears as a
      rotating current component synchronous with the rotor, producing a
      $1^{\text{st}}$-order torque ripple.
\item \textbf{Current-sensor gain mismatch.} Unequal scaling between phases
      produces a negative-sequence fundamental component and hence a
      $2^{\text{nd}}$-order torque ripple.
\item \textbf{Position-sensor error.} Resolver eccentricity or amplitude
      imbalance produces $1^{\text{st}}$- and $2^{\text{nd}}$-order errors in
      the estimated angle, which modulate the torque angle and thus the torque.
      Encoder quantisation and observer lag act similarly at higher orders.
\item \textbf{DC-link voltage ripple.} A ripple at $2f_{grid}$ (single-phase
      front end) or $6f_{grid}$ (six-pulse rectifier) propagates through the
      modulator into the currents; in voltage saturation it appears directly as
      torque ripple at the same frequency, asynchronous with the rotor.
\item \textbf{Finite current-loop bandwidth.} The current controller can only
      reject disturbances well below its crossover. At high speed $6f_{e}$ may
      exceed the achievable bandwidth, so the ripple grows with speed even
      when the machine itself is unchanged.
\end{itemize}

\subsection{Manufacturing asymmetries}

Static or dynamic rotor eccentricity, magnet dimensional and remanence
tolerances, and partial demagnetisation break the symmetry between poles and
introduce sub-harmonics, typically $1^{\text{st}}$- and $2^{\text{nd}}$-order
components, together with unbalanced magnetic pull and additional acoustic
noise.

\section{Summary of harmonic orders}

\begin{table}[h]
\centering
\small
\begin{tabular}{lll}
\toprule
Source & Order (electrical) & Current dependence \\
\midrule
Cogging & $k\,N_{c}/p$, \; $N_{c}=\operatorname{lcm}(Q,2p)$ & none \\
Back-EMF harmonics $6n\pm1$ & $6n$ & linear \\
Current harmonics $6n\pm1$ & $6n$ & linear \\
Saliency, slotting inductance modulation & $6n$ & quadratic \\
Saturation and cross-saturation & $6n$ & quadratic, load dependent \\
Dead time & $6n$ & strongest at low current \\
Sensor offset & $1$ & constant absolute value \\
Sensor gain mismatch & $2$ & linear \\
Position error (resolver) & $1$, $2$ & linear \\
Rotor eccentricity, magnet tolerance & $1$, $2$ & linear \\
DC-link ripple & asynchronous ($2f_{g}$, $6f_{g}$) & linear \\
\bottomrule
\end{tabular}
\caption{Dominant torque-ripple mechanisms and their harmonic signature.}
\end{table}

A practical consequence: the harmonic order observed in a torque or speed
spectrum identifies the mechanism. A pure $6^{\text{th}}$ order that scales with
current points to back-EMF or current distortion; a $6^{\text{th}}$ order that
scales with current squared points to saturation; a component present at zero
current is cogging; $1^{\text{st}}$- and $2^{\text{nd}}$-order components point
to the sensing chain or to mechanical asymmetry.

\section{Propagation into the DC link and into the grid}
\label{sec:grid}

This section analyses the following scenario: a back-to-back structure with an
active front end (AFE) and a motor-side inverter sharing a \emph{light} DC link
(film capacitor only, no electrolytic bank, no DC choke), feeding a PMSM at a
fundamental frequency $f_{e}=\SI{15.6}{\hertz}$ on a $f_{g}=\SI{50}{\hertz}$
grid, with \emph{no} torque-ripple compensation on the motor side. The ripple
mechanisms of Sections \ref{sec:cogging}--\ref{sec:current} are therefore
present at full amplitude, and the question is where the associated pulsating
power ends up.

\subsection{Instantaneous power balance}

Neglecting losses and the rate of change of stored magnetic energy, the power
delivered by the motor-side inverter equals the shaft power. A torque ripple
component of order $n$ therefore corresponds to a pulsating power
\begin{equation}
\tilde{p}_{n}(t) = \tilde{T}_{n}\,\omega_{m}\cos(n\omega_{e}t+\varphi_{n}),
\qquad \omega_{m}=\frac{\omega_{e}}{p},
\label{eq:pripple}
\end{equation}
so that, in per-unit terms, the pulsating power fraction equals the torque
ripple fraction:
\begin{equation}
\frac{\hat{\tilde{p}}_{n}}{P_{sh}}=\frac{\tilde{T}_{n}}{T} .
\end{equation}

\paragraph{Cogging is an exception.} Cogging torque is not supplied
electrically: it is an energy exchange between the magnet field and the rotor
inertia. It reaches the DC link only indirectly, through the speed ripple it
produces, $\Delta\hat{\omega}_{m}\approx \tilde{T}_{cog}/(J\,n\omega_{e})$, and
through the reaction of the speed and current loops. At $n\omega_{e}$ well above
the mechanical bandwidth this path is strongly attenuated by the inertia. The
DC-side pulsation is therefore dominated by the \emph{current-dependent} ripple:
back-EMF harmonics, saturation-induced inductance modulation, dead time and
sensor errors.

\subsection{Split between capacitor and grid}

The pulsating power \eqref{eq:pripple} can only be absorbed by the DC-link
capacitor or transferred to the grid through the AFE. The split is set by the
DC-voltage control loop. Denoting by $T_{v}(s)$ its complementary sensitivity
(closed-loop transfer from power disturbance to grid-side power),
\begin{equation}
\tilde{p}_{grid}=|T_{v}(jn\omega_{e})|\,\tilde{p}_{n},
\qquad
\tilde{p}_{C}=|1-T_{v}(jn\omega_{e})|\,\tilde{p}_{n},
\end{equation}
and the resulting DC-link voltage ripple is
\begin{equation}
\Delta\hat{V}_{dc,n}
=\frac{|1-T_{v}(jn\omega_{e})|\;\hat{\tilde{p}}_{n}}{n\omega_{e}\,C_{dc}\,V_{dc}}
=\frac{|1-T_{v}(jn\omega_{e})|\;\tilde{T}_{n}}{n\,p\,C_{dc}\,V_{dc}} .
\label{eq:dcripple}
\end{equation}

Equation \eqref{eq:dcripple} contains the first non-obvious result: with the
DC-voltage loop inactive at the ripple frequency, the DC-link voltage ripple
caused by torque ripple is \textbf{independent of speed}. The pulsating power
grows linearly with $\omega_{e}$, while the capacitor impedance falls as
$1/(n\omega_{e})$, and the two effects cancel exactly. Operating at low
fundamental frequency does not reduce the DC-link voltage ripple; it only
reduces the pulsating power in absolute terms.

The second result concerns the split itself. A DC-voltage loop with closed-loop
bandwidth $f_{bw}$ and a single dominant pole rejects the disturbance with
$|T_{v}|\approx f_{bw}/(nf_{e})$ above crossover. At $f_{e}=\SI{50}{\hertz}$ the
sixth-order ripple sits at \SI{300}{\hertz} and only a few percent reaches the
grid; at $f_{e}=\SI{15.6}{\hertz}$ it sits at \SI{93.6}{\hertz}, i.e.\ only about
three times the typical DC-loop bandwidth, and roughly a third of the pulsating
power is transferred to the grid. \textbf{Low-speed operation is the worst case
for grid-side emission}, even though the absolute pulsating power is smaller.

An equivalent energy view is useful for sizing: the energy that must swing in
the capacitor over one ripple period is
$\Delta W = \hat{\tilde{p}}_{n}/(n\omega_{e})$, to be compared with the stored
energy $W_{C}=\tfrac{1}{2}C_{dc}V_{dc}^{2}$, giving
$\Delta \hat{V}_{dc}/V_{dc}\approx \Delta W/(2W_{C})$.

\subsection{Worked example}

Table \ref{tab:case} collects a representative case. The torque ripple is
assumed to be \SI{3}{\percent} of rated torque at the sixth order, which is a
realistic figure for an uncompensated drive with dead time and a moderately
distorted back-EMF.

\begin{table}[h]
\centering
\small
\begin{tabular}{lll}
\toprule
Quantity & Symbol & Value \\
\midrule
Grid & $V_{g}$, $f_{g}$ & \SI{690}{\volt}, \SI{50}{\hertz} \\
Short-circuit ratio & $S_{sc}/S_{n}$ & 20 \\
Rated power & $P_{n}$ & \SI{250}{\kilo\watt} \\
DC-link voltage & $V_{dc}$ & \SI{1100}{\volt} \\
DC-link capacitance & $C_{dc}$ & \SI{200}{\micro\farad} ($\approx\SI{0.8}{\micro\farad\per\kilo\watt}$) \\
Pole pairs & $p$ & 4 \\
Fundamental frequency & $f_{e}$ & \SI{15.6}{\hertz} ($n=234$\,rpm) \\
Rated torque & $T_{n}$ & \SI{10.2}{\kilo\newton\metre} \\
\midrule
Torque ripple, 6th order & $\tilde{T}_{6}$ & \SI{3}{\percent} $=\SI{306}{\newton\metre}$ \\
Ripple frequency & $6f_{e}$ & \SI{93.6}{\hertz} \\
Pulsating power (peak) & $\hat{\tilde{p}}_{6}$ & \SI{7.5}{\kilo\watt} \\
DC current ripple (peak) & $\hat{\tilde{i}}_{dc}$ & \SI{6.8}{\ampere} \\
DC-loop bandwidth & $f_{bw}$ & \SI{30}{\hertz} $\Rightarrow |T_{v}|\approx0.32$ \\
Power to the grid & $\tilde{p}_{grid}$ & \SI{2.4}{\kilo\watt} \\
DC-link voltage ripple & $\Delta\hat{V}_{dc}$ & \SI{39}{\volt} $=\SI{3.6}{\percent}$ \\
Grid fundamental current & $I_{1}$ & \SI{209}{\ampere} \\
Each interharmonic sideband & $I_{ih}$ & \SI{1.0}{\ampere} $=\SI{0.48}{\percent}$ of $I_{1}$ \\
\bottomrule
\end{tabular}
\caption{Representative uncompensated case, AFE\,+\,inverter with light DC link.}
\label{tab:case}
\end{table}

The sideband amplitude follows from the pulsating power on the grid side. With
amplitude-invariant Park and unity displacement factor,
$\hat{\tilde{p}}_{grid}=\tfrac{3}{2}\hat{V}_{ph}\,\hat{\imath}_{d,r}$, so
\begin{equation}
\hat{\imath}_{d,r}=\frac{2\,\hat{\tilde{p}}_{grid}}{3\,\hat{V}_{ph}},
\end{equation}
and a $d$-axis component oscillating at $\omega_{r}=n\omega_{e}$ maps into a
\emph{pair} of phase-current components of amplitude $\hat{\imath}_{d,r}/2$ each.

\subsection{Interharmonic sidebands}

Because the DC-link ripple modulates the AFE, every DC-side ripple frequency
$f_{r}$ appears in the grid current as a pair of sidebands
\begin{equation}
f_{ih}^{\pm} = |f_{g} \pm f_{r}| = |f_{g} \pm n f_{e}| .
\label{eq:sidebands}
\end{equation}
Since $f_{e}$ and $f_{g}$ are unrelated, these are \emph{interharmonics}, not
harmonics. When $f_{g}-f_{r}<0$ the lower component appears with reversed phase
sequence, i.e.\ as a negative-sequence interharmonic. Table \ref{tab:ih} applies
\eqref{eq:sidebands} to the case at hand.

\begin{table}[h]
\centering
\small
\footnotesize
\begin{tabular}{lllll}
\toprule
Ripple order $n$ & Origin & $f_{r}$ [Hz] & Sidebands [Hz] & Beat with $f_{g}$ [Hz] \\
\midrule
1  & sensor offset, eccentricity & 15.6  & 34.4, 65.6   & 15.6 \\
2  & sensor gain mismatch        & 31.2  & 18.8, 81.2   & 31.2 \\
6  & back-EMF, dead time, sat. & 93.6  & \textbf{43.6} (neg.\ seq.), 143.6 & \textbf{6.4} \\
12 & back-EMF, cogging (typ.)    & 187.2 & 137.2, 237.2 & 87.2 \\
18 & higher orders               & 280.8 & 230.8, 330.8 & 180.8 \\
\bottomrule
\end{tabular}
\caption{Interharmonic emission produced by an uncompensated drive at
$f_{e}=\SI{15.6}{\hertz}$ on a \SI{50}{\hertz} grid.}
\label{tab:ih}
\end{table}

\subsection{Flicker mechanism: the $|2f_{g}-6f_{e}|$ beat}
\label{sec:flicker}

The interharmonic pair of Table \ref{tab:ih} does not merely sit near the
fundamental: it produces a genuine low-frequency fluctuation of the per-phase
power, and therefore of the per-phase voltage drop across the supply impedance.
Three equivalent derivations lead to the same beat frequency.

\paragraph{(i) Envelope view.} The lower sideband at
$f_{ih}^{-}=|f_{g}-6f_{e}|=\SI{43.6}{\hertz}$ superimposes on the
\SI{50}{\hertz} fundamental; in each phase the two components beat, and the
envelope of the phase voltage fluctuates at $|f_{g}-f_{ih}^{-}|=\SI{6.4}{\hertz}$.

\paragraph{(ii) Power view.} Taking one phase, with
$v_{a}=\hat{V}\cos(\omega_{g}t)$ and the interharmonic current
$i_{ih}=\hat{I}_{ih}\cos(\omega_{ih}t+\phi)$, the instantaneous single-phase
power is
\begin{equation}
p_{a}(t)=\frac{\hat{V}\hat{I}_{ih}}{2}
\Big[\cos\big((\omega_{g}-\omega_{ih})t-\phi\big)
+\cos\big((\omega_{g}+\omega_{ih})t+\phi\big)\Big].
\label{eq:pbeat}
\end{equation}
Substituting $\omega_{ih}=6\omega_{e}-\omega_{g}$ gives the two terms at
\begin{equation}
|2f_{g}-6f_{e}| = \SI{6.4}{\hertz}
\qquad\text{and}\qquad
6f_{e} = \SI{93.6}{\hertz},
\end{equation}
which is exactly the $\SI{100}{\hertz}-\SI{93.6}{\hertz}$ arithmetic. The two
views are algebraically identical:
\begin{equation}
f_{beat}=\Big|f_{g}-\big|f_{g}-n f_{e}\big|\Big|
= \big|2f_{g}-n f_{e}\big| \quad\text{for } n f_{e}>f_{g}.
\label{eq:fbeat}
\end{equation}
The phase of \eqref{eq:pbeat} determines how the fluctuation splits between
active and reactive power. Since the resulting voltage fluctuation at the PCC is
\begin{equation}
\frac{\Delta V}{V}\approx\frac{R\,\Delta P+X\,\Delta Q}{V^{2}},
\end{equation}
on MV networks with $X/R>3$ it is the \emph{reactive} fluctuation that dominates
the voltage modulation, and hence the flicker.

\paragraph{(iii) Flickermeter view.} Block 2 of the IEC 61000-4-15 flickermeter
squares the input voltage. The square of the sum of a \SI{50}{\hertz}
fundamental and a \SI{43.6}{\hertz} interharmonic contains a component at
\SI{6.4}{\hertz} and one at \SI{93.6}{\hertz}; the weighting filter chain, a
bandpass of roughly \SIrange{0.05}{35}{\hertz} followed by the eye--lamp
weighting, passes the first and rejects the second. The instrument therefore
responds to precisely the term identified in \eqref{eq:pbeat}.

\paragraph{Why a three-phase measurement hides it.} The lower sideband is
negative-sequence, so the \SI{6.4}{\hertz} terms of \eqref{eq:pbeat} appear in
the three phases displaced by $\mp\SI{120}{\degree}$ and cancel in the
three-phase sum: the total instantaneous power, and equivalently the space-vector
magnitude or any $dq$-domain quantity, oscillates only at \SI{93.6}{\hertz}. The
voltage drop across the supply impedance is however a per-phase quantity, and
lamps are single-phase loads. A $dq$-based analysis of the AFE will therefore
never reveal this flicker mechanism; it must be evaluated per phase.

\paragraph{Critical fundamental frequencies.} Applying \eqref{eq:fbeat} with
$n=6$, and taking the flicker-relevant band as roughly
\SIrange{0.5}{25}{\hertz} with maximum sensitivity near \SI{8.8}{\hertz}, the
sixth-order torque ripple generates flicker whenever
\begin{equation}
\big|2f_{g}-6f_{e}\big| < \SI{25}{\hertz}
\qquad\Longrightarrow\qquad
\SI{12.5}{\hertz} < f_{e} < \SI{20.8}{\hertz}
\quad (f_{g}=\SI{50}{\hertz}),
\end{equation}
with the worst case at $f_{e}=(2f_{g}-8.8)/6=\SI{15.2}{\hertz}$. Table
\ref{tab:crit} maps the band.

\begin{table}[h]
\centering
\footnotesize
\begin{tabular}{lllll}
\toprule
$f_{e}$ [Hz] & $6f_{e}$ [Hz] & $f_{ih}^{-}$ [Hz] & $f_{beat}$ [Hz] & Remark \\
\midrule
12.5  & 75.0  & 25.0 & 25.0 & lower edge of the flicker band \\
14.0  & 84.0  & 34.0 & 16.0 & inside the band \\
15.2  & 91.2  & 41.2 & 8.8  & \textbf{maximum flicker sensitivity} \\
15.6  & 93.6  & 43.6 & 6.4  & \textbf{present operating point} \\
16.67 & 100.0 & 50.0 & 0    & steady negative sequence: unbalance, not flicker \\
18.0  & 108.0 & 58.0 & 8.0  & inside the band \\
20.8  & 124.8 & 74.8 & 24.8 & upper edge of the flicker band \\
\bottomrule
\end{tabular}
\caption{Beat frequency produced by the sixth-order torque ripple as a function
of the motor fundamental frequency, $f_{g}=\SI{50}{\hertz}$.}
\label{tab:crit}
\end{table}

The operating point of \SI{15.6}{\hertz} lies \SI{0.4}{\hertz} from the
worst possible fundamental frequency for flicker generated by the sixth-order
torque ripple. This is a coincidence of the application, not a property of the
machine, and it is the single most important conclusion of this section: the
same drive at \SI{25}{\hertz} ($f_{beat}=\SI{50}{\hertz}$) or at
\SI{10}{\hertz} ($f_{beat}=\SI{40}{\hertz}$) would produce the same DC-link
ripple and the same interharmonic current, with no flicker at all.

\paragraph{Second-order mechanism: intermodulation with a real 100 Hz ripple.}
If the grid voltage is unbalanced, or the AFE current sensing is unbalanced, the
front end already produces a genuine power and DC-link ripple at
$2f_{g}=\SI{100}{\hertz}$. This component coexists with the \SI{93.6}{\hertz}
component from the motor. The superposition is linear in the capacitor voltage,
but the modulators multiply their duty cycles by $V_{dc}$, so the two ripples
\emph{intermodulate}: a component at $|2f_{g}-6f_{e}|=\SI{6.4}{\hertz}$ appears
directly in the applied voltages, producing a \emph{symmetric} sideband pair at
$f_{g}\pm\SI{6.4}{\hertz}=\SI{43.6}{}$ and \SI{56.4}{\hertz}. A symmetric pair
is pure amplitude modulation of the fundamental and is the most efficient
possible flicker source. Its amplitude is proportional to the \emph{product} of
the two relative DC-link ripples, hence second order and normally smaller than
the mechanism above, but it grows as $1/C_{dc}^{2}$ and therefore becomes
relevant exactly in the light-DC-link case considered here. A fast, unfiltered
$V_{dc}$ feedforward in both modulators removes this path; a slow or heavily
filtered one leaves it open.

\subsection{Why 15.6 Hz is an unfavourable operating point}

Three independent reasons make this particular fundamental frequency critical.

\paragraph{(a) The beat falls in the flicker band.} As shown in Section
\ref{sec:flicker}, the sixth-order ripple produces a per-phase power and voltage
fluctuation at $|2f_{g}-6f_{e}|=\SI{6.4}{\hertz}$, essentially at the minimum of
the IEC 61000-4-15 perceptibility curve. Writing $a=V_{ih}/V_{1}$, the relative
voltage fluctuation seen by the flickermeter is $\Delta V/V \approx 2a$, and
\begin{equation}
a \approx \frac{I_{ih}}{I_{n}}\cdot\frac{S_{n}}{S_{sc}}\cdot\frac{f_{ih}}{f_{g}}
\end{equation}
for a predominantly inductive grid impedance. For Table \ref{tab:case}:
$a\approx\SI{0.021}{\percent}$, hence $\Delta V/V\approx\SI{0.042}{\percent}$ and,
against a threshold of roughly \SI{0.33}{\percent} for $P_{st}=1$ at that beat
frequency, $P_{st}$ of the order of \num{0.1}. The scaling law is what matters
in practice:
\begin{equation}
P_{st}\;\propto\;\frac{\tilde{T}_{6}}{T_{n}}\cdot\frac{S_{n}}{S_{sc}} ,
\end{equation}
so the same drive on a weak grid ($S_{sc}/S_{n}=5$) with \SI{6}{\percent} torque
ripple reaches $P_{st}\approx1$, i.e.\ the planning level. Note that this
mechanism is absent on a \SI{60}{\hertz} grid at the same speed, where the
lower sideband would fall at \SI{33.6}{\hertz} and beat at \SI{26.4}{\hertz},
well outside the sensitive band.

\paragraph{(b) Proximity to the unbalance-rejection frequency.} At
$f_{e}=f_{g}/3=\SI{16.67}{\hertz}$ the sixth-order ripple would sit exactly at
$2f_{g}=\SI{100}{\hertz}$, and the lower sideband would land exactly on
\SI{50}{\hertz} with negative sequence, i.e.\ it would be indistinguishable from
grid voltage unbalance and would be actively fought by the AFE negative-sequence
controller. The operating point of \SI{15.6}{\hertz} is only \SI{1.1}{\hertz}
away from this singular condition, so the emission falls in the transition band
of the notch filters or DSC/DDSRF prefilters that are usually tuned at
\SI{100}{\hertz}: they neither pass it cleanly nor reject it, and the resulting
partial, phase-shifted response can degrade the AFE current-loop damping. A
second singular point is $f_{e}=f_{g}/6=\SI{8.33}{\hertz}$, where the lower
sideband collapses to DC and the drive injects a DC current component into the
grid, with the associated transformer-saturation risk.

\paragraph{(c) Low ripple frequency relative to the DC-voltage loop.} As shown
above, at \SI{93.6}{\hertz} the DC-voltage loop is not yet fully rolled off, so a
substantial fraction of the pulsating power is deliberately pushed onto the
grid by the controller itself.

\subsection{Secondary effects of a light DC link}

\begin{itemize}
\item \textbf{Loss of voltage headroom.} A \SI{3.6}{\percent} DC-link ripple
      reduces the voltage available to the motor-side modulator at every trough.
      If the drive operates near the modulation limit, the inverter enters
      overmodulation periodically at \SI{93.6}{\hertz}, generating additional
      low-order current harmonics which in turn increase the torque ripple: a
      positive-feedback path that does not exist with a stiff DC link.
\item \textbf{Direct modulation path.} Even with a very slow DC-voltage loop,
      the ripple reaches the grid current through the modulator: both converters
      normalise their duty cycles by $V_{dc}$, and any delay or filtering in
      that feedforward leaves a residual voltage error at \SI{93.6}{\hertz}
      which the current loops must reject. This path is independent of $T_{v}$
      and becomes dominant as $C_{dc}$ is reduced.
\item \textbf{Constant-power-load interaction.} At frequencies below its
      bandwidth the motor drive behaves as a constant-power load with negative
      incremental resistance $-V_{dc}^{2}/P$. With a light DC link the AFE
      closed-loop output impedance is larger and less damped in the
      \SIrange{50}{200}{\hertz} region, exactly where the \SI{93.6}{\hertz}
      excitation lies; the ripple can therefore be amplified rather than
      attenuated, and the impedance-ratio stability margin should be verified at
      that frequency rather than assumed.
\item \textbf{Capacitor stress.} The RMS current in a film capacitor of
      \SI{200}{\micro\farad} is dominated by the switching-frequency component,
      but the \SI{93.6}{\hertz} component adds a low-frequency thermal cycle;
      more importantly, the voltage ripple must be included in the derating and
      lifetime calculation.
\item \textbf{Measurement.} IEC 61000-4-7 uses \SI{200}{\milli\second} windows
      with \SI{5}{\hertz} resolution. The \SI{43.6}{\hertz} component falls
      between spectral lines and close to the fundamental, so leakage from the
      \SI{50}{\hertz} component can mask it; interharmonic subgroup evaluation
      and windowing must be used, otherwise the emission is simply not seen in a
      standard harmonic report.
\end{itemize}

\subsection{Summary of the propagation chain}

\begin{equation*}
\tilde{T}_{6}\;\rightarrow\;\tilde{p}=\tilde{T}_{6}\omega_{m}\;\rightarrow\;
\Big[\text{capacitor: }\Delta V_{dc}\Big]\;\Big\|\;\Big[\text{AFE: }\tilde{p}_{grid}\Big]
\;\rightarrow\; f_{g}\pm 6f_{e} \;\rightarrow\; \text{flicker, interharmonics}
\end{equation*}

With a light DC link and no motor-side compensation there is no way to make the
pulsation disappear: the design choice is only \emph{where to put it}. Either it
is absorbed by the capacitor, and then the DC-link voltage ripple and the
associated headroom loss must be budgeted, or it is passed to the grid, and then
the interharmonic emission and flicker must be assessed against IEC 61000-2-2 /
IEC 61000-3-x compatibility levels, EN 61800-3 and, where applicable,
IEEE 519. Reducing the ripple at the source on the motor side is the only option
that avoids both.

\section{Mitigation}

\subsection{Machine design}

\begin{itemize}
\item Slot/pole combination with large $\operatorname{lcm}(Q,2p)$ and small
      $\gcd(Q,2p)$.
\item Continuous or step skewing of stator slots or rotor magnets.
\item Magnet shaping: pole-arc ratio optimisation, eccentric or sinusoidally
      profiled magnet surface, magnet segmentation and pole shifting.
\item Semi-closed slots, reduced slot opening, magnetic slot wedges, dummy
      slots or tooth notches to increase the effective slotting frequency.
\item Winding layout (fractional-slot concentrated versus distributed) chosen
      considering that it changes both the back-EMF harmonic content and the
      cogging periodicity.
\end{itemize}

Skewing reduces cogging and back-EMF harmonics but slightly reduces the
fundamental torque and complicates manufacturing; the design is a compromise
between ripple, torque density and cost.

\subsection{Control}

Control-side compensation acts on the current reference or on the voltage
command:

\begin{itemize}
\item \textbf{Harmonic current injection (feedforward).} The reference is
      augmented with the harmonics that cancel the measured ripple:
      \begin{equation}
      i_{q}^{*}(\thetae)=i_{q0}^{*}+\sum_{n}A_{n}\cos(6n\thetae+\phi_{n}),
      \end{equation}
      with $A_{n}$, $\phi_{n}$ identified offline and stored in a
      speed/load-indexed lookup table.
\item \textbf{Resonant (PR / vector resonant) controllers} added in parallel to
      the synchronous-frame PI at the ripple frequencies,
      \begin{equation}
      G_{r}(s)=\frac{2k_{r}\,\omega_{c}\,s}{s^{2}+2\omega_{c}s+(6n\omega_{e})^{2}},
      \end{equation}
      providing high loop gain at $6n\omega_{e}$ without affecting the DC
      behaviour.
\item \textbf{Repetitive control} in the position domain, whose internal model
      covers all harmonics of the mechanical or electrical period at once; it is
      the natural choice when the disturbance is position-periodic (cogging)
      rather than time-periodic.
\item \textbf{Iterative learning control} for repetitive operating cycles.
\item \textbf{Disturbance observers} estimating the periodic load torque from
      the speed signal, with the estimate fed forward to the current reference.
\item \textbf{Dead-time and voltage-error compensation}, plus current-sensor
      offset and gain self-calibration, which remove the $1^{\text{st}}$,
      $2^{\text{nd}}$ and part of the $6^{\text{th}}$ order without any
      additional dynamics.
\end{itemize}

Two limits apply to all active methods. First, they require loop bandwidth at
$6n\omega_{e}$, which becomes the binding constraint at high speed; multiple
synchronous frames or position-domain formulations relieve, but do not remove,
this constraint. Second, cancelling torque ripple means deliberately injecting
current harmonics: this costs additional copper and iron losses and voltage
headroom, and the benefit must be weighed against those costs, particularly in
the flux-weakening region.

\subsection{Grid-side and DC-link measures}

When the motor-side ripple cannot be removed, the emission can still be shaped:

\begin{itemize}
\item \textbf{Notch in the DC-voltage loop} at $6f_{e}$ (and $12f_{e}$), tuned
      adaptively on the measured motor frequency. The AFE then presents a high
      impedance at the ripple frequency and stops re-exporting the pulsation to
      the grid; the price is a larger DC-link voltage ripple, which must fit in
      the modulation headroom of both converters.
\item \textbf{Deliberate capacitor sizing.} From \eqref{eq:dcripple}, the
      capacitance needed for a given ripple budget is
      $C_{dc}\geq \tilde{T}_{n}/(n\,p\,V_{dc}\,\Delta \hat{V}_{dc})$, which is
      independent of speed and therefore a hard sizing constraint rather than a
      worst-speed calculation.
\item \textbf{Active power decoupling} on a small auxiliary branch, which
      buffers the $6f_{e}$ energy without oversizing the main DC-link capacitor.
\item \textbf{Speed operating-point selection.} From \eqref{eq:fbeat}, keeping
      $f_{e}$ outside \SIrange{12.5}{20.8}{\hertz} (for $f_{g}=\SI{50}{\hertz}$)
      moves the sixth-order beat out of the flicker band entirely, and avoiding
      $f_{e}$ near $f_{g}/6$ and $f_{g}/3$ removes the DC-injection and
      apparent-unbalance conditions. Where the process allows a different gear
      ratio or a different pole count, this is by far the cheapest measure: it
      costs nothing and removes the dominant mechanism rather than attenuating
      it.
\item \textbf{DC-link voltage feedforward} with adequate bandwidth in both
      modulators, so that the ripple does not translate into an applied-voltage
      error.
\item \textbf{Impedance-based stability check} of the AFE output impedance
      against the motor-drive input impedance in the \SIrange{50}{200}{\hertz}
      range, which is the frequency band excited by the torque ripple at this
      operating point.
\end{itemize}

\end{document}
